\documentclass[11pt]{article}

\usepackage[margin=1in]{geometry}
\usepackage{amsmath,amssymb,amsthm,mathtools,bm}
\usepackage{booktabs,array,multirow}
\usepackage{graphicx}
\usepackage{microtype}
\usepackage{xcolor}
\usepackage[hypertexnames=false]{hyperref}
\usepackage{enumitem}
\usepackage{float}

\hypersetup{
  colorlinks=true,
  linkcolor=blue!55!black,
  citecolor=blue!55!black,
  urlcolor=blue!55!black
}

\newtheorem{theorem}{Theorem}
\newtheorem{proposition}[theorem]{Proposition}
\newtheorem{lemma}[theorem]{Lemma}
\newtheorem{corollary}[theorem]{Corollary}
\newtheorem{assumption}{Assumption}
\theoremstyle{definition}
\newtheorem{definition}{Definition}
\newtheorem{remark}{Remark}

\newcommand{\R}{\mathbb R}
\newcommand{\E}{\mathbb E}
\newcommand{\Pp}{\mathbb P}
\newcommand{\Sph}{\mathbb S}
\newcommand{\Id}{\mathbf I}
\newcommand{\uvec}{\mathbf u}
\newcommand{\svec}{\mathbf s}
\newcommand{\vvec}{\mathbf v}
\newcommand{\Avec}{\mathbf A}
\newcommand{\Hvec}{\mathbf H}
\newcommand{\Pvec}{\mathbf P}
\newcommand{\xivec}{\boldsymbol\xi}
\newcommand{\PiProd}{\boldsymbol\Pi}
\newcommand{\ind}{\mathbf 1}

\title{\textbf{Kesten Tails and Nonlinear Self-Stabilization in}\\
Finite-Step Random-Covariate SGD}
\author{Anonymous draft}
\date{\today}

\begin{document}
\maketitle

\begin{abstract}
We study constant-step stochastic gradient descent as an exact discrete
random dynamical system.  For batch-quadratic linear regression, the
linearized iterates form an affine random recurrence.  Below its top
Lyapunov boundary, multivariate Kesten theory predicts power-law stationary
tails.  Above that boundary the linear recursion is not stationary, but a
local negative quartic and positive higher-order correction can stabilize a
flip mode through the exact two-step map.  We define the Lyapunov exponent
and an inverse-moment ``harmonic'' diagnostic from full products of the batch
matrices and, under explicit spectral-identification, differentiability, and
strict-convexity assumptions on the pressure, prove that any inverse-moment
crossing cannot precede the Lyapunov crossing.  We also give an exact
two-dimensional deterministic flip theorem.  In the post-Lyapunov regime we
prove, for an exact scalar regenerative excursion, the discrete
on--off-intermittency law
$\Pp(|a|\leq r)\sim C r^\alpha$ and
$p(a)\sim C'|a|^{\alpha-1}$, where
$\E M^{-\alpha}=1$.  A separate proposition gives the conditional matrix
extension under projective spectral-gap, Markov-renewal, angular, and
regeneration hypotheses, with $k_\eta(-\alpha)=1$.  The analysis makes
explicit that fixed additive noise restricts the law to an intermediate
window rather than the literal origin.
We then test the nonlinear stochastic recursion directly in dimensions two and
five, with fresh random covariates, noncommuting batch Hessians, deterministic
and bounded random quartic coefficients, and calibrated step-size sweeps across the
entire post-Lyapunov, pre-harmonic interval.  Every reported trajectory uses
the finite-step SGD update itself; no iterate is clipped or reset.  The
experiments distinguish a broad centered cloud from directional bimodality
and measure how nonlinear strength, coefficient variability, temporal
persistence, batch size, and covariance anisotropy affect bursts and modes.
An additional five-noise-scale audit shows that the primary-noise chains
yield no admissible window under our stated operational criterion; the
product-root exponent emerges only after sufficient low-noise scale
separation and is stable across the tested finite horizons.
\end{abstract}

\section{Introduction}

Let $B_t$ denote an independently sampled minibatch and consider the exact
finite-step recursion
\begin{equation}
  \uvec_{t+1}
  =\Avec_{B_t}\uvec_t+\eta\xivec_{B_t}
   -\eta\nabla V_{B_t}(\uvec_t),
  \qquad
  \Avec_B=\Id-\eta\Hvec_B.
  \label{eq:intro-recursion}
\end{equation}
The product
$\PiProd_n=\Avec_{B_{n-1}}\cdots\Avec_{B_0}$ governs the local linear
dynamics.  When its top Lyapunov exponent is negative, the affine part of
\eqref{eq:intro-recursion} can possess a stationary law with a Kesten tail
\cite{Kesten1973,Goldie1991}.  When the exponent is positive, the affine model
alone cannot be stationary.  The central question is then whether the
nonlinear term returns expanding trajectories, and whether its stationary
cloud remains centered or separates into modes.

We keep three statements distinct throughout:
\begin{enumerate}[leftmargin=2em]
  \item a negative top Lyapunov exponent is a criterion for linear
  stationarity;
  \item recurrence of the nonlinear chain is a separate global property;
  \item an inverse-moment crossing is a local diagnostic and is not, without
  additional structure, a universal necessary-and-sufficient condition for
  bimodality.
\end{enumerate}

The main contributions are as follows.
\begin{enumerate}[label=\arabic*.,leftmargin=2em]
  \item We specify a batch-quadratic loss with an exact, smooth,
  globally controlled negative-quartic continuation and derive its ordinary
  finite-step SGD update.
  \item We show why the untamed explicit polynomial update has an escape
  region, even when the highest-order term is restoring in the loss.
  \item We define the full-product Lyapunov and pressure diagnostics and give
  an interacting-particle estimator for the inverse-moment pressure.
  \item For an exact scalar regenerative coordinate, we prove the
  on--off-intermittency law
  $\Pp(|a|\leq r)\sim C r^\alpha$ and
  $p(a)\sim C'|a|^{\alpha-1}$, where
  $\E M^{-\alpha}=1$.  We give a separate conditional Markov-additive
  extension with $k_\eta(-\alpha)=1$ and isolate the additional projective
  spectral-gap, regeneration, and cone conditions needed for matrix products.
  \item For an explicit separable two-dimensional flip normal form, we prove
  that one symmetric period-two orbit is attracting and a second, outer orbit
  is repelling under explicit coefficient conditions.
  \item We perform calibrated finite-step experiments below the Lyapunov edge,
  throughout the Lyapunov--harmonic gap, and beyond the harmonic diagnostic.
  Deterministic, i.i.d., and persistent random quartic coefficients are tested
  at fixed mean.
  \item We separate a broad centered cloud from two-lobe structure by comparing
  broad-variance and pinned-direction covariates in dimensions two and five.
\end{enumerate}

\section{Related work}

\paragraph{Random affine recursions and heavy-tailed SGD.}
Kesten's renewal theorem and Goldie's implicit renewal theory are the classical
sources for power-law tails of affine stochastic recursions
\cite{Kesten1973,Goldie1991}; the modern multivariate theory is synthesized by
Buraczewski, Damek, and Mikosch~\cite{BuraczewskiDamekMikosch2016}.  In
stochastic optimization, G\"urb\"uzbalaban, \c{S}im\c{s}ekli, and Zhu prove
heavy-tailed stationary behavior for constant-step SGD on quadratic linear
regression~\cite{GurbuzbalabanSimsekliZhu2021}, while Hodgkinson and Mahoney
identify multiplicative noise as a general mechanism for heavy tails in
discrete stochastic optimizers~\cite{HodgkinsonMahoney2021}.  Our
sub-Lyapunov theorem is a specialization of this line.  The present focus is
what can remain recurrent after the affine recursion has lost stationarity,
when an explicit finite-amplitude nonlinear feedback is retained.

\paragraph{Edge of stability and self-stabilization.}
Cohen et al. documented that full-batch neural-network training commonly
operates near the classical flip boundary~\cite{CohenEtAl2021}.  Subsequent
work analyzed unstable convergence and edge-of-stability dynamics from
complementary perspectives~\cite{AhnZhangSra2022,AroraLiPanigrahi2022}.
Damian, Nichani, and Lee identified a deterministic cubic self-stabilization
mechanism and its implicit bias toward lower sharpness
\cite{DamianNichaniLee2023}.  We preserve their emphasis on the actual
finite-step map, but study fresh random covariates, noncommuting matrix
products, and the distinction between recurrence, shells, and directional
bimodality.

\paragraph{On--off intermittency and reinjection.}
The post-Lyapunov mechanism studied here is an instance of classical
on--off intermittency: multiplicative motion away from a transversely
unstable set is repeatedly reinjected by nonlinear dynamics.  This mechanism
and its bursting and laminar-phase statistics were introduced and developed
by Platt, Spiegel, and Tresser; Heagy, Platt, and Hammel; and Ott and
Sommerer~\cite{PlattSpiegelTresser1993,HeagyPlattHammel1994,OttSommerer1994}.
We therefore do not claim the small-amplitude power law as a new
intermittency mechanism.  In particular, the classical $-3/2$ law for
laminar durations is a first-passage statement and is not the stationary
amplitude exponent proved below.  Our role is to formulate the Cram\'er--renewal
argument for an exact scalar regenerative recursion, give a conditional
Markov-additive extension tied to the full matrix-product pressure, and state
explicitly the reinjection, projective-mixing, and observable-window
conditions.  The proof uses standard
random-walk renewal and exponential-tilting tools
\cite{Asmussen2003,NeyNummelin1987I,NeyNummelin1987II}.

\paragraph{Random dynamical systems and bifurcation language.}
The separation between sample-path stability and changes in an invariant
distribution is standard in the random-dynamical-systems viewpoint
\cite{Arnold1998}.  Here the Lyapunov crossing concerns the inner linear
cocycle, whereas a modal change is a property of the nonlinear invariant law.
Keeping those statements separate prevents the inverse-moment diagnostic
from being misidentified as a universal bimodality theorem.

\section{Exact finite-step stochastic quadratic--quartic SGD}
\label{sec:model}

\subsection{Random-covariate regression}

For a minibatch of size $b$, let
$\mathbf X_B\in\R^{b\times d}$ contain the covariates and let
$\boldsymbol\varepsilon_B\in\R^b$ contain the label innovations.  Define
\begin{equation}
  \Hvec_B=\frac1b\mathbf X_B^\top\mathbf X_B,
  \qquad
  \xivec_B=\frac1b\mathbf X_B^\top\boldsymbol\varepsilon_B.
  \label{eq:hessian-noise}
\end{equation}
The batch loss is
\begin{equation}
  L_B(\uvec)
  =\frac1{2b}\|\mathbf X_B\uvec-\boldsymbol\varepsilon_B\|^2
   +V_{\beta_B}(\uvec).
  \label{eq:batch-loss}
\end{equation}

Let $\Pvec$ be the orthogonal projection onto the critical subspace and put
$s=\|\Pvec\uvec\|^2$.  For $R>0$ and $\beta>0$, define the exact saturating
potential
\begin{equation}
  V_\beta(\uvec)
  =-\frac{\beta R^2}{2}
      \left[s-R^2\log\!\left(1+\frac{s}{R^2}\right)\right].
  \label{eq:saturating-potential}
\end{equation}
It is a genuine smooth loss, not an update rule imposed after the SGD step.
Its gradient is
\begin{equation}
 \nabla V_\beta(\uvec)
 =-\frac{\beta R^2s}{R^2+s}\Pvec\uvec.
 \label{eq:saturating-gradient}
\end{equation}
Near the reference point,
\begin{equation}
  V_\beta(\uvec)
  =-\frac\beta4\|\Pvec\uvec\|^4
   +\frac{\beta}{6R^2}\|\Pvec\uvec\|^6
   +O(\|\Pvec\uvec\|^8).
  \label{eq:local-expansion}
\end{equation}
Thus the loss has the requested negative quartic and positive sextic terms
locally.  As $\|\Pvec\uvec\|\to\infty$, its gradient is asymptotic to
$-\beta R^2\Pvec\uvec$.  The outer discrete matrix cocycle therefore supplies
a rigorous sufficient condition for global recurrence; we do not claim that
this condition is necessary.  Section~\ref{sec:recurrence} gives the precise
condition used below.

Substituting \eqref{eq:saturating-gradient} into the gradient step gives the
exact recursion used in every experiment:
\begin{equation}
  \uvec_{t+1}
  =(\Id-\eta\Hvec_{B_t})\uvec_t+\eta\xivec_{B_t}
   +\eta\frac{\beta_tR^2s_t}{R^2+s_t}\Pvec\uvec_t,
  \quad s_t=\|\Pvec\uvec_t\|^2.
  \label{eq:exact-discrete-update}
\end{equation}
Since $\nabla^2V_\beta(\mathbf0)=\mathbf0$, changing the nonlinear
coefficient process does not change the linear multiplier
$\Avec_B=\Id-\eta\Hvec_B$.

\subsection{Constant, i.i.d., and persistent coefficients}

We compare three bounded processes with the same mean $\bar\beta$:
\begin{align}
  &\text{constant:} &&\beta_t=\bar\beta,\notag\\
  &\text{i.i.d.:} &&\beta_t=\bar\beta(1+\delta S_t),
      \quad S_t\stackrel{\mathrm{i.i.d.}}{\sim}
      \operatorname{Unif}\{-1,+1\},\notag\\
  &\text{persistent:} &&\beta_t=\bar\beta(1+\delta S_t),
      \quad \Pp(S_{t+1}=-S_t)=p_{\mathrm{flip}}.
  \label{eq:beta-processes}
\end{align}
Here $0\leq\delta<1$, so $\beta_t$ is positive and bounded.  The persistent case is a homogeneous
Markov chain on the augmented state $(\uvec_t,S_t)$.  Thus slowly changing
curvature is represented directly at the iteration level.  The experiment
uses $p_{\rm flip}=0.002$ and $0.00035$, corresponding to stay probabilities
$0.998$ and $0.99965$.

\subsection{Why the raw polynomial cannot be used globally}

The Taylor polynomial of \eqref{eq:saturating-gradient} is useful locally,
but its explicit gradient step is not a valid global stabilizer.

\begin{proposition}[Uniform escape region for an explicit polynomial map]
\label{prop:polynomial-escape}
Let
\begin{equation}
  x_{t+1}=a_tx_t+g_tx_t^3-hx_t^5+e_t,
  \qquad h>0,
  \label{eq:polynomial-map}
\end{equation}
where $|a_t|\leq A$, $|g_t|\leq G$, and $|e_t|\leq E$ for every $t$.
There exists a finite $R_0=R_0(A,G,E,h)$ such that
\begin{equation}
  |x_t|\geq R_0
  \quad\Longrightarrow\quad
  |x_{t+1}|\geq\frac h2|x_t|^5\geq2|x_t|.
  \label{eq:escape-bound}
\end{equation}
Consequently, every trajectory that enters this region diverges in norm.
\end{proposition}

\begin{proof}
The reverse triangle inequality gives
\[
 |x_{t+1}|
 \geq h|x_t|^5-G|x_t|^3-A|x_t|-E.
\]
Choose $R_0$ so that the last three terms are at most
$h|x_t|^5/2$ and $h|x_t|^5/2\geq2|x_t|$ whenever
$|x_t|\geq R_0$.  Iterating the inequality proves divergence.
\end{proof}

\begin{remark}
A positive highest-order coefficient in a polynomial loss does not prevent an
explicit fixed-step update from overshooting.  Projection, clipping, and
resetting would conceal this issue.  Equation~\eqref{eq:saturating-potential}
instead changes the loss itself outside the local regime and remains an
ordinary SGD experiment.
\end{remark}

\subsection{A sufficient discrete recurrence condition}
\label{sec:recurrence}

The far-field multiplier of \eqref{eq:exact-discrete-update} is
\begin{equation}
  \Avec_{\infty,t}
  =\Id-\eta\Hvec_t+\eta\beta_tR^2\Pvec.
  \label{eq:outer-multiplier}
\end{equation}
The primary experiments estimate its top product exponent.  The theorem below
uses the stronger $m$-step block-moment condition.

\begin{theorem}[Discrete self-stabilization by an outer cocycle]
\label{thm:foster}
Suppose $(\Hvec_t,\xivec_t,\beta_t)$ is i.i.d.,
$0\leq\beta_t\leq\beta_{\max}$, and, for some $p\in(0,1]$ and integer
$m\geq1$,
\begin{equation}
  \E\left\|
    \Avec_{\infty,m-1}\cdots\Avec_{\infty,0}
  \right\|^p\leq\rho<1,
  \qquad
  \E\|\Avec_{\infty,t}\|^p<\infty,
  \qquad
  \E\|\xivec_t\|^p<\infty.
  \label{eq:far-field-moment}
\end{equation}
Then the exact update \eqref{eq:exact-discrete-update} satisfies
\begin{equation}
  \E[1+\|\uvec_{t+m}\|^p\mid\uvec_t=\uvec]
  \leq \rho(1+\|\uvec\|^p)+C
  \label{eq:foster-drift}
\end{equation}
for some $\rho<1$ and $C<\infty$.  If the chain is Feller, the drift condition
implies existence of an invariant probability law with a finite $p$th
moment.  If, in addition,
the chain is irreducible and aperiodic and compact sets are petite, the
invariant law is unique and the chain is positive Harris recurrent.  Under
all the preceding drift and kernel assumptions, a positive inner exponent
$\gamma_\eta>0$ together with \eqref{eq:far-field-moment} gives an exact
finite-step self-stabilization regime: typical infinitesimal perturbations of
the inner cocycle grow, whereas the nonlinear chain is positive Harris
recurrent.
\end{theorem}

\begin{proof}[Proof sketch]
Write the exact update as
$\uvec_{t+1}=\Avec_{\infty,t}\uvec_t+\mathbf b_t(\uvec_t)
+\eta\xivec_t$, where
\[
 \mathbf b_t(\uvec)
 =-\eta\beta_t\frac{R^4}{R^2+\|\Pvec\uvec\|^2}\Pvec\uvec.
\]
Since $\sup_{r\geq0}r/(R^2+r^2)=1/(2R)$,
$\|\mathbf b_t(\uvec)\|\leq\eta\beta_{\max}R^3/2$ uniformly.  Iterate for
$m$ steps and use $\|\sum_j\mathbf y_j\|^p\leq\sum_j\|\mathbf y_j\|^p$
for $p\leq1$.  Independence and \eqref{eq:far-field-moment} yield
\eqref{eq:foster-drift}; standard $m$-step Foster theory gives the invariant
law and the Harris conclusion.
\end{proof}

The experiments below estimate the full-product outer Lyapunov exponent
\begin{equation}
 \gamma_{\infty}
 =\lim_{n\to\infty}\frac1n
 \log\|\Avec_{\infty,n-1}\cdots\Avec_{\infty,0}\|.
 \label{eq:outer-lyapunov}
\end{equation}
They do not by themselves numerically certify
\eqref{eq:far-field-moment}.  Under pressure regularity and moments in a
neighborhood of zero, $\gamma_{\infty}<0$ implies that such a block condition
holds for some sufficiently small $p>0$ and sufficiently large $m$.
Accordingly, the recurrence theorem is applied conditional on these standard
pressure assumptions.

For the experimental transition kernel, the Feller property is immediate:
the update is continuous in $\uvec$ for every realization of the bounded
driving variables, so dominated convergence maps bounded continuous test
functions to continuous functions.  We have not proved
$\psi$-irreducibility, aperiodicity, or petiteness of compact sets for this
specific law.  This is nontrivial because the residual signs are Rademacher;
any smoothing must come from the continuously distributed Hessian parameters
$(\lambda,\theta)$.  Thus uniqueness and positive Harris recurrence in
Theorem~\ref{thm:foster} remain conditional on its stated kernel assumptions.

\section{Discrete random-matrix diagnostics}
\label{sec:matrix-diagnostics}

\subsection{Lyapunov exponent and projective pressure}

For the linear multiplier product
\begin{equation}
  \PiProd_n=\Avec_{B_{n-1}}\cdots\Avec_{B_0},
  \qquad \Avec_B=\Id-\eta\Hvec_B,
  \label{eq:product}
\end{equation}
define the top Lyapunov exponent
\begin{equation}
  \gamma_\eta
  =\lim_{n\to\infty}\frac1n\log\|\PiProd_n\|
  \quad\text{almost surely}.
  \label{eq:lyapunov}
\end{equation}
For $q$ in a range where the required moments are finite, define the
projective transfer operator
\begin{equation}
  (\mathcal P_{\eta,q}f)(\svec)
  =\E\left[
    \|\Avec_B\svec\|^q
    f\left(\frac{\Avec_B\svec}{\|\Avec_B\svec\|}\right)
  \right],
  \qquad \svec\in\Sph^{d-1},
  \label{eq:transfer}
\end{equation}
and
\begin{equation}
  k_\eta(q)=\rho(\mathcal P_{\eta,q}),
  \qquad
  \Lambda_\eta(q)=\log k_\eta(q).
  \label{eq:pressure}
\end{equation}

\begin{assumption}[Negative-pressure spectral identification]
\label{ass:negative-pressure-identification}
For $q\in[-1,0]$, the operator $\mathcal P_{\eta,q}$ has a simple isolated
leading eigenvalue on a suitable H\"older space and, for every
$\svec\in\Sph^{d-1}$,
\begin{equation}
 \log\rho(\mathcal P_{\eta,q})
 =\lim_{n\to\infty}\frac1n
   \log\E\|\PiProd_n\svec\|^q,
 \label{eq:negative-pressure-identification}
\end{equation}
with a value independent of $\svec$.  The pressure is differentiable at
$q=0$ and satisfies $\Lambda_\eta'(0)=\gamma_\eta$.
\end{assumption}

\begin{remark}[Convexity from finite products]
For fixed $n$ and $\svec$, the map
$q\mapsto n^{-1}\log\E\|\PiProd_n\svec\|^q$ is convex by H\"older's
inequality.  Whenever the finite limit in
\eqref{eq:negative-pressure-identification} exists, its pointwise limit is
therefore convex.  The substantive unverified ingredient for the experimental
matrix law is the spectral identification, equivalently an appropriate
spectral-gap statement, especially at $q=-1$.  Strict convexity remains a
separate nondegeneracy assumption.
\end{remark}

Under the standard spectral hypotheses,
$\Lambda_\eta(0)=0$ and
$\Lambda_\eta'(0)=\gamma_\eta$
\cite{GuivarchLePage2012,WangSteinsaltz2020}.

When $\gamma_\eta<0$ and the affine recurrence is nondegenerate, a positive
root
\begin{equation}
  k_\eta(\kappa_\eta)=1,
  \qquad \kappa_\eta>0,
  \label{eq:kesten-root}
\end{equation}
is the Kesten tail exponent under the usual irreducibility, moment, and
non-arithmetic assumptions.  This statement concerns the linear affine
recursion.  For the nonlinear chain, the Kesten fit is expected only before
the nonlinear force is appreciable.

\begin{theorem}[Multivariate Kesten baseline]
\label{thm:kesten-baseline}
Consider the affine recursion
$\uvec_{t+1}=\Avec_t\uvec_t+\eta\xivec_t$ with i.i.d.
$(\Avec_t,\xivec_t)$.  Assume $\gamma_\eta<0$, the required logarithmic and
$\kappa$-moments are finite, the matrix semigroup is strongly irreducible and
proximal, the associated product is non-arithmetic, and the recursion has no
common deterministic fixed point.  If $k_\eta(\kappa_\eta)=1$ for some
$\kappa_\eta>0$, then the affine recursion has a stationary law $\mathbf U$ that is
multivariately regularly varying with index $\kappa_\eta$.  In particular,
\begin{equation}
 \Pp(\|\mathbf U\|>t)\sim C t^{-\kappa_\eta},
 \qquad
 \Pp(|\vvec^\top\mathbf U|>t)\sim C(\vvec)t^{-\kappa_\eta},
 \qquad t\to\infty,
 \label{eq:kesten-directional}
\end{equation}
for every direction charged by the spectral tail measure.  The exponent is
common; only $C(\vvec)$ depends on direction.
\end{theorem}

Theorem~\ref{thm:kesten-baseline} is the quadratic reference result tested in
the sub-Lyapunov control.  It is not extended unchanged beyond
$\gamma_\eta=0$, where the affine stationary law ceases to exist.

\begin{definition}[Two finite-step diagnostics]
Along a step-size path, the \emph{Lyapunov edge} is a simple crossing
$\gamma_\eta=0$.  Whenever the negative moment is finite, define
\begin{equation}
  h_\eta=k_\eta(-1)
  \label{eq:harmonic}
\end{equation}
and call $h_\eta=1$ the \emph{inverse-moment harmonic diagnostic}.
\end{definition}

\begin{theorem}[Ordering of the two diagnostics]
\label{thm:ordering}
Suppose Assumption~\ref{ass:negative-pressure-identification} holds on
$[-1,0]$.
If $h_\eta<1$, then $\gamma_\eta>0$.  At a crossing $h_\eta=1$, strict
convexity on $[-1,0]$ implies $\gamma_\eta>0$.
\end{theorem}

\begin{proof}
Convexity and $\Lambda_\eta(0)=0$ imply
\[
  \gamma_\eta=\Lambda_\eta'(0)
  \geq\Lambda_\eta(0)-\Lambda_\eta(-1)
  =-\log h_\eta.
\]
This proves the conclusion when $h_\eta<1$.
For the strict crossing case, $\Lambda_\eta(-1)=\Lambda_\eta(0)=0$ and
strict convexity imply $\Lambda_\eta(q)<0$ for every $q\in(-1,0)$.  Hence
\[
 \gamma_\eta=\Lambda_\eta'(0)
 \geq\frac{\Lambda_\eta(0)-\Lambda_\eta(q)}{-q}>0.
\]
\end{proof}

Theorem~\ref{thm:ordering} is an ordering result, not a theorem that the
second crossing creates modes.  If $\gamma_\eta>0$ and a unique negative root
exists, write
\begin{equation}
  k_\eta(-\alpha_\eta)=1,
  \qquad \alpha_\eta>0.
  \label{eq:negative-root}
\end{equation}
Assume additionally that $-1$ and $-\alpha_\eta$ lie in the finite domain of
$\Lambda_\eta$, that $-\alpha_\eta$ is its unique negative zero, and that
$\Lambda_\eta$ is strictly convex on the convex hull of
$\{0,-1,-\alpha_\eta\}$.  Then the algebraic equivalences are
\begin{equation}
  h_\eta>1\Longleftrightarrow0<\alpha_\eta<1,
  \qquad
  h_\eta=1\Longleftrightarrow\alpha_\eta=1,
  \qquad
  h_\eta<1\Longleftrightarrow\alpha_\eta>1.
  \label{eq:root-order}
\end{equation}

\subsection{A proved discrete reinjection law}
\label{sec:reinjection-theorem}

The exponent in \eqref{eq:negative-root} describes a \emph{small-amplitude}
law after the Lyapunov crossing, not a Kesten upper tail.  The mechanism is
classical on--off intermittency.  In logarithmic amplitude, typical motion
has positive drift away from the origin, while rare contracting products
create deep negative excursions; nonlinear self-stabilization returns the
process from the outer scale and starts another excursion.  The following
renewal-theoretic theorem makes that statement exact for a scalar
regenerative coordinate.  A direction-dependent matrix product is handled
separately in Proposition~\ref{prop:matrix-reinjection}.  The scalar theorem
is a discrete random-walk result and uses no diffusion approximation.
For a scalar flip coordinate, $M_n$ below is the absolute one-step
multiplier; the alternating sign is carried separately and does not affect
the log amplitude.

\begin{theorem}[Scalar regenerative reinjection law]
\label{thm:scalar-reinjection}
Let $(M_n)_{n\geq1}$ be i.i.d. positive random variables, put
$Y_n=\log M_n$ and $S_n=\sum_{j=1}^nY_j$, and assume:
\begin{enumerate}[label=(\roman*),leftmargin=2.2em]
  \item $0<\mu:=\E Y_1<\infty$, and the law of $Y_1$ is non-arithmetic;
  \item for some $\alpha>0$ and $\epsilon>0$,
  \begin{equation}
     \E e^{-\alpha Y_1}=1,
     \qquad \E e^{-(\alpha+\epsilon)Y_1}<\infty,
     \qquad
     m_\alpha:=-\E[Y_1e^{-\alpha Y_1}]\in(0,\infty);
     \label{eq:scalar-cramer-assumptions}
  \end{equation}
  \item the reinjection variable $Z<0$ is independent of $(Y_n)$ and
  $\E e^{-(\alpha+\epsilon)Z}<\infty$;
  \item for $X_n=Z+S_n$ and
  $\tau=\inf\{n\geq0:X_n\geq0\}$, one has $\E\tau<\infty$.
\end{enumerate}
Construct a regenerative chain by retaining
$(X_0,\ldots,X_{\tau-1})$, discarding the overshoot, drawing an independent
copy of $Z$, and repeating.  Its invariant cycle-occupation law $\pi_X$
satisfies, for every fixed $h>0$,
\begin{equation}
 \pi_X((x,x+h])
 \sim
 \frac{D_\alpha}{\E\tau}
 \frac{e^{\alpha h}-1}{\alpha m_\alpha}e^{\alpha x},
 \qquad x\to-\infty,
 \label{eq:log-reinjection-local-tail}
\end{equation}
where
\begin{equation}
 D_\alpha
 :=\E\!\left[e^{-\alpha Z}-e^{-\alpha X_\tau}\right]>0.
 \label{eq:scalar-reinjection-constant}
\end{equation}
If the tilted renewal measure admits a density for which the renewal-density
theorem holds, then $\pi_X$ has a left-tail density $q_X$ and
\begin{equation}
 q_X(x)\sim
 \frac{D_\alpha}{m_\alpha\E\tau}e^{\alpha x}.
 \label{eq:log-reinjection-density}
\end{equation}
Consequently, for $R=e^X$,
\begin{equation}
 \Pp(R\leq r)\sim
 \frac{D_\alpha}{\alpha m_\alpha\E\tau}r^\alpha,
 \qquad
 p_R(r)\sim
 \frac{D_\alpha}{m_\alpha\E\tau}r^{\alpha-1},
 \qquad r\downarrow0.
 \label{eq:radial-reinjection-density}
\end{equation}
For a reflection-symmetric signed coordinate $A$ with $|A|=R$,
\begin{equation}
 p_A(a)\sim
 \frac{D_\alpha}{2m_\alpha\E\tau}|a|^{\alpha-1},
 \qquad a\to0^\pm.
 \label{eq:signed-reinjection-density}
\end{equation}
\end{theorem}

\begin{proof}
Let $F$ be the law of $Y_1$ and introduce the Cram\'er tilt
$F_\alpha(dy)=e^{-\alpha y}F(dy)$.  It is a probability law by
\eqref{eq:scalar-cramer-assumptions}, and its mean is
$\int yF_\alpha(dy)=-m_\alpha<0$.  If
$U=\sum_{n\geq0}F^{*n}$ and
$U_\alpha=\sum_{n\geq0}F_\alpha^{*n}$ are the corresponding potential
measures, then the change of measure gives the exact identity
\begin{equation}
 U(dt)=e^{\alpha t}U_\alpha(dt).
 \label{eq:renewal-tilt-identity}
\end{equation}
The non-arithmetic random-walk renewal theorem, applied to the tilted walk
of drift $-m_\alpha$, therefore yields for a walk started at $z$,
\begin{equation}
 U_z((x,x+h])
 \sim e^{-\alpha z}
 \frac{e^{\alpha h}-1}{\alpha m_\alpha}e^{\alpha x},
 \qquad x\to-\infty.
 \label{eq:started-renewal-asymptotic}
\end{equation}
The $(\alpha+\epsilon)$ moment of $Z$ justifies averaging this limit over
the reinjection law.

Let
$G_Z(B)=\E\sum_{n=0}^{\tau-1}\ind_{\{X_n\in B\}}$ be the occupation
measure of one killed excursion.  The strong Markov property at $\tau$
gives the exact killed-potential decomposition
\begin{equation}
 U_Z(B)=G_Z(B)+\E[U_{X_\tau}(B)].
 \label{eq:killed-potential-identity}
\end{equation}
Using \eqref{eq:started-renewal-asymptotic} in the two potential terms
shows that their leading coefficients differ by $D_\alpha$.  Notice that
$D_\alpha>0$ pointwise because $Z<0\leq X_\tau$.  Finally, the
renewal--reward formula gives $\pi_X=G_Z/\E\tau$, proving
\eqref{eq:log-reinjection-local-tail}.  Summing fixed-width log intervals,
with the standard exponential renewal bound controlling the remainder,
gives the small-ball CDF in \eqref{eq:radial-reinjection-density}.  The
density version of the renewal theorem gives
\eqref{eq:log-reinjection-density}; the changes of variables
$p_R(r)=q_X(\log r)/r$ and $R=|A|$ give
\eqref{eq:radial-reinjection-density}--\eqref{eq:signed-reinjection-density}.
This is the usual Cram\'er--renewal calculation
\cite{Asmussen2003,Alsmeyer1994}.
\end{proof}

\begin{corollary}[What the Lyapunov--harmonic gap means]
\label{cor:scalar-reinjection-order}
Under the assumptions of Theorem~\ref{thm:scalar-reinjection}, including its
renewal-density hypothesis, let
$\psi(s)=\log\E M_1^{-s}$ be finite and strictly convex on
$[0,s_*]$ for some $s_*\geq1$.
If $\E\log M_1>0$ and $\E M_1^{-1}>1$, then the nonzero root
$\psi(\alpha)=0$ is unique and satisfies $0<\alpha<1$.  Hence
\eqref{eq:signed-reinjection-density} is an integrable but divergent central
cusp.  At $\E M_1^{-1}=1$, the nonzero root is $\alpha=1$ and the leading
local density is flat.  If $\E M_1^{-1}<1$ and a nonzero root exists in
$(1,s_*]$, then that root has $\alpha>1$ and the idealized local density
vanishes at the origin.  This last conclusion is central depletion, not by
itself bimodality.
\end{corollary}

\begin{proof}
Since $\psi(0)=0$ and
$\psi'(0)=-\E\log M_1<0$, strict convexity together with
$\psi(1)>0$ produces exactly one second zero in $(0,1)$.  The equality and
post-crossing cases follow from the same convexity argument on the domain
where the second root exists.
\end{proof}

Thus the apparent paradox in the gap is resolved: the origin is unstable in
the typical Lyapunov sense, yet rare contracting products create so much
occupation near it that the stationary density diverges there.  The nonlinear
self-stabilization term fixes the reinjection law, the normalization, and the
outer cutoff; because its derivative vanishes at the origin, it does not
change the Cram\'er exponent of the inner multiplicative walk.

The same separation explains bounded stochastic self-stabilization
coefficients.  An i.i.d. $\beta_t$, or a persistent finite-state $\beta_t$
after augmenting the regenerative state, can change $Z$, $D_\alpha$,
$\E\tau$, and the cutoff distribution.  It does not change the inner
pressure root so long as the nonlinear derivative at the origin remains zero
and the linear matrix law is unchanged.  This is a conditional invariance of
the local exponent, not a claim that the full stationary distribution is
independent of the coefficient process.

\begin{remark}[Additive core and nonlinear cutoff]
\label{rem:reinjection-window}
Theorem~\ref{thm:scalar-reinjection} is an ideal excursion theorem.  At fixed
nonzero additive noise, its $r\downarrow0$ conclusion cannot hold literally:
the additive core rounds the density.  If a family of exact two-boundary
excursions has lower radius $r_{\rm core}$, upper radius $r_{\rm nl}$,
reinjection inside the two boundaries, convergence in the tilted
$(\alpha+\epsilon)$ moment, and the uniform local-renewal bound needed for a
triangular-array renewal theorem, the killed-potential proof gives the same
exponent at radii $r_L$ satisfying
\begin{equation}
 \frac{r_{\rm core}}{r_L}\to0,
 \qquad
 \frac{r_L}{r_{\rm nl}}\to0.
 \label{eq:explicit-intermediate-window}
\end{equation}
Thus the operational implication is the two-sided window
\begin{equation}
 r_{\rm core}\ll r\ll r_{\rm nl}.
 \label{eq:reinjection-window-short}
\end{equation}
The required uniform renewal and reinjection conditions are additional
hypotheses, not consequences of a large scale ratio.
\end{remark}

For an actual additive recursion, the lower boundary is an approximation,
not a literal reset.  A useful deterministic error bound makes the required
scale separation explicit.  If
$a_{n+1}=B_{n+1}a_n+N_{n+1}(a_n)+\sigma\zeta_{n+1}$ with
$|B_n|\geq b_*>0$, $|N_n(a)|\leq C|a|^{1+\delta}$, and
$|\zeta_n|\leq K$, then on $r_-\leq|a_n|\leq r_+$,
\begin{equation}
 \left|
  \log\frac{|a_{n+1}|}{|a_n|}-\log|B_{n+1}|
 \right|
 \leq \frac{2}{b_*}
 \left(Cr_+^\delta+\frac{\sigma K}{r_-}\right),
 \label{eq:local-log-error-bound}
\end{equation}
provided the right-hand side is at most one.  This follows by factoring
$B_{n+1}a_n$ and using $|\log|1+z||\leq2|z|$ for $|z|\leq1/2$.
Hence a transfer from the exact excursion theorem to nonlinear SGD requires
both $\sigma\ll r_-\ll r_+\ll r_{\rm nl}$ and a verified regenerative
minorization for the nonlinear return kernel.  We do not infer either fact
from a fitted slope.

\begin{proposition}[Conditional matrix and projection extension]
\label{prop:matrix-reinjection}
Let $(\Avec_n)$ be i.i.d. invertible matrices.  On real projective space
$\mathbb P^{d-1}$, define the Markov additive process
\begin{equation}
 J_n=[\Avec_n\svec_{n-1}],
 \qquad
 X_n=X_0+\sum_{j=1}^n\log\|\Avec_j\svec_{j-1}\|.
 \label{eq:projective-additive-process}
\end{equation}
Here $[\svec]$ denotes the line through either unit representative
$\svec$, and $\svec_{n-1}$ is any unit representative of $J_{n-1}$.
Suppose the induced projective chain is geometrically ergodic on a
H\"older space; the induced operator $\mathcal P_{\eta,q}$ has a simple
isolated leading eigenvalue in a neighborhood of $q=-\alpha$;
$k_\eta(-\alpha)=1$ and
$m_\alpha:=-\Lambda_\eta'(-\alpha)>0$; and the Markov additive process is
non-arithmetic and satisfies the moment hypotheses of the Markov renewal
theorem.  Let $r_\alpha>0$ be its leading eigenfunction.  At every upper
exit, suppose the nonlinear return restarts $(X_0,J_0)\sim\nu$ independently
of future matrices, and, for
$\tau=\inf\{n\geq0:X_n\geq0\}$, assume
\begin{equation}
 \E_\nu\tau<\infty,
 \quad
 \int e^{-(\alpha+\epsilon)x}r_\alpha(j)\,\nu(dx,dj)<\infty,
 \quad
 D_\alpha^{\rm mat}:=E_\nu\!\left[
 e^{-\alpha X_0}r_\alpha(J_0)
 -e^{-\alpha X_\tau}r_\alpha(J_\tau)
 \right]>0.
 \label{eq:matrix-regeneration-assumptions}
\end{equation}
Then there exist $C_\alpha\in(0,\infty)$ and a probability measure
$\chi_\alpha$ on $\mathbb P^{d-1}$ such that, for every fixed $h>0$ and
every $\chi_\alpha$-continuity set $D$,
\begin{equation}
 \lim_{x\to-\infty}e^{-\alpha x}
 \pi\!\left(X\in(x,x+h],\ J\in D\right)
 =C_\alpha\frac{e^{\alpha h}-1}{\alpha}\chi_\alpha(D).
 \label{eq:matrix-log-reinjection}
\end{equation}
Writing $\mathbf U=e^X\svec$ with either projective representative, this
implies the radial small-ball law
\begin{equation}
 \Pp(\|\mathbf U\|\leq r)
 \sim\frac{C_\alpha}{\alpha}r^\alpha.
 \label{eq:matrix-radial-small-ball}
\end{equation}
For a fixed direction $\vvec$, suppose the excursion law is pinned away from
its orthogonal hyperplane,
\begin{equation}
 |\vvec^\top\svec|\geq c_\vvec>0
 \quad \pi\text{-almost surely throughout occupied excursion states}.
 \label{eq:pinned-projection-condition}
\end{equation}
Then the absolute projection has the explicit small-ball law
\begin{equation}
 \Pp(|\vvec^\top\mathbf U|\leq r)
 \sim\frac{C_\alpha}{\alpha}r^\alpha
 \int_{\mathbb P^{d-1}}|\vvec^\top\svec|^{-\alpha}
 \chi_\alpha(d[\svec]).
 \label{eq:matrix-projection-small-ball}
\end{equation}
If the corresponding Markov-renewal density theorem holds, it yields density
asymptotics proportional to $r^{\alpha-1}$, with constants matching
\eqref{eq:matrix-radial-small-ball}--\eqref{eq:matrix-projection-small-ball}.
A signed-side statement
requires augmenting the projective state by sign or parity; reflection
symmetry then makes the two side constants equal.
\end{proposition}

\begin{proof}[Proof logic]
The Doob transform on projective space
\begin{equation}
 (\widehat{\mathcal P}_\alpha f)(j)
 =\frac1{r_\alpha(j)}
 \E\!\left[
   \|\Avec_1\svec\|^{-\alpha}
   r_\alpha([\Avec_1\svec])
   f([\Avec_1\svec])
 \right]
 \label{eq:projective-doob-transform}
\end{equation}
is a probability kernel, where $\svec$ is either representative of $j$.  On
$n$-step paths its likelihood ratio is
$e^{-\alpha(X_n-X_0)}r_\alpha(J_n)/r_\alpha(J_0)$, and the tilted
log-radius drift is
$\Lambda_\eta'(-\alpha)=-m_\alpha<0$.  The Markov renewal theorem
\cite{Kesten1974,NeyNummelin1987I,NeyNummelin1987II} therefore gives a
constant left-tail renewal measure and a limiting tilted angular law.
Undoing the transform contributes $e^{\alpha x}$.  The weighted defect in
\eqref{eq:matrix-regeneration-assumptions}, the killed-potential identity,
and renewal--reward normalization then give
\eqref{eq:matrix-log-reinjection}.  Summing log intervals gives
\eqref{eq:matrix-radial-small-ball}.  Finally,
$\log|\vvec^\top\mathbf U|=X+\log|\vvec^\top\svec|$; pinning makes the
second term bounded, so integrating the radial small-ball limit over
$\chi_\alpha$ gives \eqref{eq:matrix-projection-small-ball}.
\end{proof}

Proposition~\ref{prop:matrix-reinjection} identifies the precise additional
high-dimensional issue: the pressure root controls the radial log excursion,
but a projection also depends on direction.  Working on projective space
removes the deterministic sign flip visible in the experimental cocycle; a
signed statement instead uses a parity-augmented chain or the two-step
skeleton.  Without a projective spectral gap, an i.i.d. regenerative angular
restart, and pinning away from $\vvec^\perp$, a radial cusp need not imply the
same directional cusp.
These assumptions are not silently declared verified for the nonlinear
experimental chain.

\subsection{Full-product estimators}

The top exponent is estimated with normalized vector iteration.  Given one
matrix stream and a unit vector $\svec_0$, iterate
\begin{equation}
  \ell_t=\log\|\Avec_t\svec_t\|,
  \qquad
  \svec_{t+1}=\frac{\Avec_t\svec_t}{\|\Avec_t\svec_t\|},
  \qquad
  \widehat\gamma_\eta=\frac1T\sum_{t=0}^{T-1}\ell_t.
  \label{eq:gamma-estimator}
\end{equation}
The implementation averages over independent product trajectories and
discloses every deterministic seed used for the plotted point estimates.

For $k_\eta(q)$ we use a population Feynman--Kac estimator.  Maintain
$N_p$ directions $\svec_t^{(i)}$.  At step $t$:
\begin{enumerate}[leftmargin=2.2em]
  \item independently sample a batch matrix $\Avec_t^{(i)}$ for every
  particle;
  \item set
  $w_t^{(i)}=\|\Avec_t^{(i)}\svec_t^{(i)}\|^q$ and
  $\widetilde\svec_{t+1}^{(i)}=
  \Avec_t^{(i)}\svec_t^{(i)}/
  \|\Avec_t^{(i)}\svec_t^{(i)}\|$;
  \item record
  $c_t=N_p^{-1}\sum_iw_t^{(i)}$ and resample the new particles from
  $\widetilde\svec_{t+1}^{(i)}$ with probabilities proportional to
  $w_t^{(i)}$.
\end{enumerate}
Then
\begin{equation}
  \log\widehat k_\eta(q)
  =\frac1T\sum_{t=0}^{T-1}\log c_t.
  \label{eq:smc-pressure}
\end{equation}
The particle count, product length, burn-in, and seeds used for each plot are
recorded in the executable artifact.  For a generic matrix law, finiteness of
the corresponding inverse moment remains a theorem assumption and cannot be
proved by a finite simulation.  For the bounded laws used here,
Lemma~\ref{lem:inverse-moment-finite} proves finiteness on the full calibration
grid.  What remains assumed is the projective spectral identification in
Assumption~\ref{ass:negative-pressure-identification}.
In particular, $h_\eta$ is not replaced by an average of individual Hessian
eigenvalues or by a one-step operator norm.

\section{Exact discrete flip mechanism}
\label{sec:flip}

\subsection{Two-step local calculation}

Along a signed critical coordinate, the local update has the form
\begin{equation}
  x_{t+1}=-(1+\delta_t)x_t+g_tx_t^3-h_tx_t^5+O(x_t^7).
  \label{eq:local-flip}
\end{equation}
At exact flip and to cubic order,
$f(x)=-x+gx^3+O(x^5)$, hence
\begin{equation}
  f(f(x))=x-2gx^3+O(x^5).
  \label{eq:two-step-expansion}
\end{equation}
For the negative local quartic in \eqref{eq:local-expansion}, $g>0$.
Thus the one-step cubic points outward in the raw coordinate, but it is
restoring in the exact two-step flip map.

\subsection{An exact two-dimensional theorem}

The following theorem gives a finite-step mode-forming mechanism without any
limiting approximation.

\begin{theorem}[Exact symmetric two-cycles in two dimensions]
\label{thm:exact-2d-flip}
Let $g,h>0$, $|\rho|<1$, and
\begin{equation}
  F_\delta(x,y)
  :=\bigl(-(1+\delta)x+gx^3-hx^5,\;\rho y\bigr).
  \label{eq:2d-map}
\end{equation}
Assume
\begin{equation}
  0<\delta<\frac{g^2}{4h},
  \qquad
  \Delta=g^2-4h\delta,
  \qquad
  z_\pm=\frac{g\pm\sqrt\Delta}{2h}.
  \label{eq:cycle-roots}
\end{equation}
Then $F_\delta$ has two symmetric period-two orbits
\begin{equation}
  \mathcal O_\pm
  =\{(\sqrt{z_\pm},0),(-\sqrt{z_\pm},0)\}.
  \label{eq:two-cycles}
\end{equation}
The outer orbit $\mathcal O_+$ is unstable.  The inner orbit
$\mathcal O_-$ is locally asymptotically stable if
\begin{equation}
  z_-\sqrt\Delta<1.
  \label{eq:cycle-stability}
\end{equation}
It is unstable if $z_-\sqrt\Delta>1$.  At
$z_-\sqrt\Delta=1$ the longitudinal multiplier of the two-step map equals
one, and no conclusion is asserted without a separate higher-order
nondegeneracy analysis.
The origin is unstable.  Every trajectory in the basin of
$\mathcal O_-$ has the time-averaged limiting measure
\begin{equation}
  \frac12\delta_{(\sqrt{z_-},0)}
  +\frac12\delta_{(-\sqrt{z_-},0)}.
  \label{eq:cycle-measure}
\end{equation}
\end{theorem}

\begin{proof}
Because the map is odd in $x$, a point $(r,0)$ belongs to a symmetric
period-two orbit precisely when its first coordinate maps to $-r$.  For
$r\neq0$, this condition is
\[
  hz^2-gz+\delta=0,
  \qquad z=r^2,
\]
which gives \eqref{eq:cycle-roots}.  At either orbit the derivative of the
first-coordinate map is, after using the cycle equation,
\[
  f'(r)=-1+2gz-4hz^2.
\]
Therefore
$f'(\sqrt{z_-})=-1+2z_-\sqrt\Delta$ and
$f'(\sqrt{z_+})=-1-2z_+\sqrt\Delta$.
The first-coordinate multiplier of the two-step map is $f'(r)^2$.  For
$z_-$ it is less than one when $z_-\sqrt\Delta<1$ and greater than one when
$z_-\sqrt\Delta>1$; at equality it is one.  For $z_+$ it is strictly greater
than one.  The transverse two-step multiplier is $\rho^2<1$.
Finally, the derivative at the origin has first-coordinate magnitude
$1+\delta>1$.  Convergence to the attracting orbit implies
\eqref{eq:cycle-measure}.
\end{proof}

\begin{remark}
Theorem~\ref{thm:exact-2d-flip} applies exactly to maps orthogonally conjugate
to \eqref{eq:2d-map}.  For a general pinned critical direction with
cross-coordinate nonlinearities or noncommuting random matrices, persistence
of this orbit requires a separate hyperbolic-persistence or random
center-manifold argument and is not asserted by this theorem.  Determining
whether additive innovations broaden the two phases into two visible modes
is a property of the full invariant kernel, not of the local derivative
alone.
\end{remark}

The globally completed loss admits an equally explicit deterministic result.

\begin{corollary}[Exact two-dimensional flip for the saturated loss]
\label{cor:saturated-flip}
Let $\mathbf M\in\R^{2\times2}$ be symmetric with eigenvalues $m_1>m_2$ and leading unit
eigenvector $\vvec_1$.  For $L=\eta\beta R^2$, consider
\begin{equation}
 F(\mathbf z)=\left(\mathbf M
 -L\frac{\|\mathbf z\|^2}{R^2+\|\mathbf z\|^2}\Id\right)\mathbf z.
 \label{eq:saturated-demodulated-map}
\end{equation}
Put $d=m_1-1$.  A pair of nonzero fixed points along $\vvec_1$ exists if and
only if $0<d<L$; this pair is
\begin{equation}
 \mathbf z_\pm=\pm R\sqrt{\frac{d}{L-d}}\,\vvec_1.
 \label{eq:saturated-mode-location}
\end{equation}
Their longitudinal and transverse Jacobian multipliers are
\begin{equation}
 1-2d\left(1-\frac dL\right),
 \qquad 1+m_2-m_1,
 \label{eq:saturated-mode-multipliers}
\end{equation}
Hence the two points are locally asymptotically stable if both multipliers
in \eqref{eq:saturated-mode-multipliers} have modulus strictly less than one,
and they are unstable if at least one multiplier has modulus strictly greater
than one.  In the remaining boundary case, when both multipliers have modulus
at most one and at least one has modulus one, the point is nonhyperbolic and
the Jacobian calculation alone is inconclusive.  When
$\mathbf M=\eta\mathbf H-\mathbf I$ arises from the sign
demodulation of the SGD recurrence above, these two fixed points are the two
phases of one period-two orbit in the original coordinates.
\end{corollary}

\begin{proof}
A nonzero fixed point must be an eigenvector of $\mathbf M$.  Along
$\vvec_1$, the fixed-point equation is
$Lr^2/(R^2+r^2)=m_1-1=d$, giving
\eqref{eq:saturated-mode-location}.  Differentiating
\eqref{eq:saturated-demodulated-map} in the radial and transverse directions
gives \eqref{eq:saturated-mode-multipliers}.
\end{proof}

\subsection{Stochastic and time-varying coefficients}

The coefficient process changes finite-amplitude feedback without changing
the derivative at the origin.  The following identity measures that feedback
directly from the finite-step recurrence.

\begin{proposition}[Exact discrete nonlinear balance]
\label{prop:nonlinear-balance}
Let a noise-free stationary solution satisfy
\begin{equation}
 \mathbf Z_{t+1}
 =\left(\mathbf M_t-a_t(R_t)\Id\right)\mathbf Z_t,
 \quad R_t=\|\mathbf Z_t\|,\quad
 \mathbf S_t=\mathbf Z_t/R_t,
 \label{eq:noise-free-balance-map}
\end{equation}
where $R_t>0$ almost surely, $\E|\log R_t|<\infty$, and
$a_t(r)=\eta\beta_tR^2r^2/(R^2+r^2)$.  Then
\begin{equation}
 \E\log\left\|
   (\mathbf M_t-a_t(R_t)\Id)\mathbf S_t
 \right\|=0.
 \label{eq:nonlinear-balance}
\end{equation}
For a stationary solution of the corresponding noisy recurrence, if
$\Pp(R_t>0)=1$ and the logarithms below are integrable, the exact balance is
instead
\begin{equation}
 \E\log\frac{\|(
 \mathbf M_t-a_t(R_t)\Id)\mathbf Z_t+\eta\boldsymbol\zeta_t\|}
 {\|\mathbf Z_t\|}=0,
 \label{eq:noisy-nonlinear-balance}
\end{equation}
\end{proposition}

\begin{proof}
Both statements follow by taking expectations of
$\log R_{t+1}-\log R_t$ under stationarity.
\end{proof}

For i.i.d. $\beta_t$, the coefficient is part of the ordinary transition
kernel.  For persistent Markov $\beta_t$, stationarity concerns the augmented
chain $(\mathbf u_t,\beta_t)$.  A periodic schedule has a periodic invariant
family, not one time-independent distribution; arbitrary deterministic drift
is therefore summarized phase-conditionally in our experiments.

Theorem~\ref{thm:foster} applies directly to the i.i.d. case.  For the
persistent Markov case, the analogous rigorous conclusion requires an
$m$-step drift condition conditional on the coefficient state for the
augmented chain.  Finite simulated trajectories alone do not establish this
condition.

\begin{proposition}[Directional density-curvature identity]
\label{prop:mode-curvature}
Suppose a stationary discrete chain has the form
$\mathbf Z_{t+1}=F_{E_t}(\mathbf Z_t)+\boldsymbol\epsilon_t$.  For a unit
vector $\vvec$, suppose that, conditional on $E_t$, the projected innovation
$\vvec^\top\boldsymbol\epsilon_t$ is independent of $\mathbf Z_t$, has an
even $C^2$ density $g_{E_t,\vvec}$, and differentiation may be passed through
the stationary expectation.  Let $p_{\vvec}$ be the stationary density of
$\vvec^\top\mathbf Z$, and put
$Y=\vvec^\top F_{E_t}(\mathbf Z_t)$.  Then
\begin{equation}
 p_{\vvec}''(0)
 =\E\left[
 g_{E_t,\vvec}''(-Y)
 \right].
 \label{eq:exact-origin-curvature}
\end{equation}
If $p_{\vvec}$ is even and $p_{\vvec}''(0)\neq0$, then the origin is a strict
local minimum if and only if $p_{\vvec}''(0)>0$.  If, in addition,
$p_{\vvec}(y)\to0$ as $|y|\to\infty$ and $p_{\vvec}$ has at most one critical
point on $(0,\infty)$, then
\begin{equation}
 p_{\vvec}\text{ is bimodal}
 \quad\Longleftrightarrow\quad p_{\vvec}''(0)>0.
 \label{eq:bimodal-iff-curvature}
\end{equation}
\end{proposition}

\begin{proof}
Stationarity gives
$p_{\vvec}(y)=\E[g_{E_t,\vvec}(y-Y)]$.  Differentiate twice under the
expectation and set $y=0$.  The final equivalence follows from evenness and
the stated critical-point restriction.
\end{proof}

Proposition~\ref{prop:mode-curvature} is first an exact curvature identity.  It
becomes an if-and-only-if bimodality criterion only under the explicit
evenness, decay, nonzero-curvature, and one-positive-critical-point shape
assumptions above; the last assumption excludes competing shoulders or modes
and is substantive.  This is not a general consequence of $k_\eta(-1)=1$.
For Gaussian projected noise with variance $\tau^2$, the integrand is
$\phi_\tau(Y)(Y^2/\tau^4-1/\tau^2)$.  The primary experiments use Rademacher
residuals, whose projected innovations need not have the required $C^2$
density, so their directional densities are classified directly and
\eqref{eq:exact-origin-curvature} is not invoked.

\section{Shells versus directional bimodality}
\label{sec:shells}

An isotropic nonlinear chain need not choose a direction.  If its radius is
concentrated near $r$ while its angle is nearly uniform, each fixed
projection can remain centered.

\begin{proposition}[A shell can have centered projections]
\label{prop:shell}
If $\uvec$ is uniform on the sphere of radius $r$ in $\R^d$, then for every
unit vector $\vvec$, the density of $A=\vvec^\top\uvec$ is
\begin{equation}
  p_A(a)\propto
  (r^2-a^2)^{(d-3)/2}\ind_{\{|a|<r\}}.
  \label{eq:sphere-projection}
\end{equation}
For $d>3$, this density has its unique maximum at zero.
\end{proposition}

Thus the experiments separately report the norm, the joint angular
distribution, and projections onto both population and empirical critical
directions.  A nonzero radial mode is called a shell; it is not called
bimodality unless a directional marginal has two reproducible maxima.

\section{Finite-step experimental protocol}
\label{sec:protocol}

\subsection{Data and exact updates}

The primary design uses bounded, full-rank random covariates so both the
inner and outer discrete products are well defined.  Fix a population
direction $\vvec_\star=\mathbf e_1$.  For block $j=1,\ldots,B$, draw
\begin{equation}
 \lambda_{t,j}=1+\delta_\lambda Z_{t,j},\qquad
 Z_{t,j}\sim\operatorname{Unif}[-1,1],
 \label{eq:random-eigenvalue}
\end{equation}
and
\begin{equation}
 \vvec_{t,j}=\cos(\theta_{t,j})\mathbf e_1
 +\sin(\theta_{t,j})\mathbf q_{t,j},
 \qquad \theta_{t,j}\sim\operatorname{Unif}[-\phi,\phi],
 \label{eq:random-direction}
\end{equation}
where $\mathbf q_{t,j}$ is Haar-uniform on the unit sphere in
$\mathbf e_1^\perp$.  Define
\begin{equation}
 \Hvec_{t,j}=h_\perp\Id+
 (\lambda_{t,j}-h_\perp)\vvec_{t,j}\vvec_{t,j}^\top,
 \qquad h_\perp=0.8,
 \qquad
 \Hvec_t=\frac1B\sum_{j=1}^B\Hvec_{t,j}.
 \label{eq:bounded-random-hessian}
\end{equation}
This is exactly a regression Gram matrix.  Taking
$\mathbf X_{t,j}=\sqrt d\,\Hvec_{t,j}^{1/2}$ and stacking the $B$ blocks gives
\begin{equation}
 \frac1{dB}\mathbf X_t^\top\mathbf X_t=\Hvec_t.
 \label{eq:exact-gram-realization}
\end{equation}
For independent Rademacher residual vectors
$\boldsymbol\varepsilon_{t,j}\in\{\pm\sigma_\varepsilon\}^d$, we use the
exact regression innovation
\begin{equation}
 \xivec_t=\frac{1}{B\sqrt d}\sum_{j=1}^B
 \Hvec_{t,j}^{1/2}\boldsymbol\varepsilon_{t,j}
 =\frac1{dB}\mathbf X_t^\top\boldsymbol\varepsilon_t.
 \label{eq:bounded-regression-innovation}
\end{equation}
The placement of $B$ outside the square root is essential.  Conditional on
the Gram blocks,
\[
 \operatorname{Cov}(\xivec_t\mid\Hvec_{t,1:B})
 =\frac{\sigma_\varepsilon^2}{dB}\Hvec_t,
\]
so the innovation variance decreases as $B^{-1}$.  The implementation uses
exactly the first expression in \eqref{eq:bounded-regression-innovation}.
An independent $240{,}000$-draw audit at each $B\in\{1,2,4,8\}$ gives
empirical-to-theoretical ratios for $\E\|\xivec_t\|^2$ of
$1.0000,1.0003,0.9993,0.9978$, respectively; moreover,
$B\E\|\xivec_t\|^2\simeq2.25\times10^{-5}$ in all four cases.  Thus both the
implementation and its exported audit exhibit the intended $B^{-1}$ variance
scaling.
Thus every trajectory is generated by the ordinary finite-step gradient of
one freshly sampled regression loss.  We do not apply a secondary integrator,
sample from a target density, clip a state, or reset a run.

In all experiments in this section we take $\Pvec=\Id$ and $R=1$; the fixed
direction $\vvec_\star$ refers to the anisotropy of the random Hessian law,
not to a rank-one nonlinear projection.

Because the directions $\vvec_{t,j}$ are freshly sampled, the Hessians are
noncommuting with probability one.  We use two matrix laws.  The
\emph{broad-gap} law has $(\delta_\lambda,\phi)=(0.4,0.18)$ and is designed
to resolve a sizeable interval between the two product diagnostics.  The
\emph{clean-mode} law has $(\delta_\lambda,\phi)=(0.2,0.10)$ and a more
strongly pinned critical direction, making directional mode separation
visible at finite noise.  Both use $h_\perp=0.8$, $R=1$, and
$\sigma_\varepsilon=0.005$.

The phrase ``pinned'' has an exact local geometric meaning for the clean
linear product.  Let
$\mathcal C_\phi=\{\svec:\mathbf e_1^\top\svec\geq
\cos(\phi)\|\svec\|\}$.

\begin{lemma}[Invariant alternating cones for the clean law]
\label{lem:clean-cone}
For the clean-mode law, $\lambda_{t,j}\geq h_\perp$ almost surely.  If
$\phi<\pi/4$ and $\eta h_\perp>1$, then
\begin{equation}
 (\eta\Hvec_t-\Id)\mathcal C_\phi\subseteq\mathcal C_\phi,
 \qquad
 (\Id-\eta\Hvec_t)\mathcal C_\phi\subseteq-\mathcal C_\phi.
 \label{eq:clean-invariant-cones}
\end{equation}
Thus the noiseless linear cocycle alternates between two signed cones and
stays a fixed positive angle away from $\mathbf e_1^\perp$.
\end{lemma}

\begin{proof}
For $\svec\in\mathcal C_\phi$, every
$\vvec_{t,j}\in\mathcal C_\phi$ and
$\vvec_{t,j}^\top\svec>0$ because $2\phi<\pi/2$.  Moreover,
\begin{equation}
 (\eta\Hvec_t-\Id)\svec
 =(\eta h_\perp-1)\svec
 +\frac{\eta}{B}\sum_{j=1}^B
 (\lambda_{t,j}-h_\perp)
 (\vvec_{t,j}^\top\svec)\vvec_{t,j}.
 \label{eq:clean-cone-decomposition}
\end{equation}
All coefficients are nonnegative, and the circular cone is convex.  This
proves the first inclusion; the second follows by changing sign.
\end{proof}

Lemma~\ref{lem:clean-cone} verifies local cone preservation, but not global
entry into a cone under additive noise or the regenerative return hypothesis
of Proposition~\ref{prop:matrix-reinjection}.  It therefore supplies one
piece of direction selection rather than a complete matrix cusp theorem for
the experiment.

\begin{lemma}[Finite inverse moments on the calibration grid]
\label{lem:inverse-moment-finite}
For the broad-gap law,
$0.6\Id\preceq\Hvec_t\preceq1.4\Id$ almost surely; for the clean-mode law,
$0.8\Id\preceq\Hvec_t\preceq1.2\Id$ almost surely.  Consequently, for every
tested $\eta\in[1.82,2.42]$, every $q<0$, every $n\geq1$, and every unit
vector $\svec$,
\[
  \E\|\PiProd_n\svec\|^q<\infty.
\]
In particular, the inverse moment defining $k_\eta(-1)$ is finite throughout
the reported calibration grid.
\end{lemma}

\begin{proof}
Averaging Gram blocks preserves their spectral bounds.  If
$a\Id\preceq\Hvec_t$, then $\eta a>1$ on the stated grid and
$\sigma_{\min}(\Id-\eta\Hvec_t)\geq\eta a-1=:c_\eta>0$.  Hence
$\|\PiProd_n\svec\|\geq c_\eta^n$, and for $q<0$,
$\E\|\PiProd_n\svec\|^q\leq c_\eta^{nq}<\infty$.
\end{proof}

\subsection{Threshold calibration}

We estimate $\gamma_\eta$ and $h_\eta$ for the two-dimensional broad-gap law
at $B\in\{1,2,4,8\}$ and for the two-dimensional clean-mode law at $B=1$.
Let
$\eta_{\mathrm L}$ denote the interpolated zero of $\widehat\gamma_\eta$ and
$\eta_{\mathrm H}$ the interpolated zero of
$\log\widehat h_\eta$ when a crossing is observed.
Lemma~\ref{lem:inverse-moment-finite} guarantees inverse-moment finiteness on
this grid.  Interpreting the particle estimate as
$\rho(\mathcal P_{\eta,-1})$ additionally uses
Assumption~\ref{ass:negative-pressure-identification}, which is not proved
for the specific simulated law.  The broad-gap distribution experiment uses
one point below the Lyapunov edge, three points spanning the gap, and two
points beyond the harmonic edge.  A separate scaling experiment uses five evenly spaced
relative positions $r\in\{0.1,0.3,0.5,0.7,0.9\}$, where
\begin{equation}
 \eta(r)=\eta_{\mathrm L}
 +r(\eta_{\mathrm H}-\eta_{\mathrm L}).
 \label{eq:relative-grid}
\end{equation}
The clean-mode experiment uses one point below $\eta_{\mathrm L}$, one in
the gap, and two beyond $\eta_{\mathrm H}$.  The five-dimensional check uses
a step size selected from the two-dimensional clean-mode calibration; we
also calibrate its product edges independently to verify, rather than assume,
its position relative to the five-dimensional harmonic edge.

Threshold uncertainty is Monte Carlo uncertainty conditional on the local
linear interpolation used to locate each zero.  For $\eta_{\mathrm L}$ we
resample complete, independent product trajectories.  For
$\eta_{\mathrm H}$ we rerun the full Feynman--Kac particle system with
independent seeds and bootstrap those independent SMC replicates.  The
reported 95\% intervals use $600$ product trajectories, $18$ independent SMC
replicates, and $3{,}000$ bootstrap resamples.  They do not include grid or
model-specification uncertainty.

\subsection{Stationarity and distributional measurements}

Each primary configuration advances many independent chains under the exact
recursion, discards a fixed burn-in, and stores every tenth iterate.  The
exact chain counts and horizons are reported with the final run.  The
reliability experiment uses additional independent seeds and a
finite-radius initialization check.  We inspect the following diagnostics:
\begin{enumerate}[leftmargin=2em]
  \item the median, $0.9$, and $0.99$ quantiles of $\|\uvec_t\|$;
  \item histograms of the fixed projection $\vvec_\star^\top\uvec_t$;
  \item central occupancy, exceedance fraction, and lag-one log-radius
  correlation;
  \item nonlinear tangent growth and the exact conditional log-radius
  balance.
\end{enumerate}
All primary runs must complete without a nonfinite iterate.  Untamed controls
are stopped at first escape and reported separately.

For each stored iterate, define the saturation fraction
\begin{equation}
  a_t=\frac{\|\uvec_t\|^2}{1+\|\uvec_t\|^2}.
  \label{eq:activation-ratio}
\end{equation}
Its empirical mean records how strongly the bounded nonlinear feedback is
activated.  We also report the exact conditional log-radius decomposition
and the tangent exponent of the full nonlinear update.

Mode claims use 301 equal-width bins followed by Gaussian smoothing with
standard deviation $1.5$ bins.  On the clean-mode range $[-1.1,1.1]$ this is
a coordinate bandwidth of $0.01096$.  We report the ratio of the smoothed
density at zero to its positive-side maximum.  Its uncertainty is computed by
a chain-block bootstrap that resamples complete chains rather than correlated
retained iterates.  A projection is called split only when the upper endpoint
of its 95\% bootstrap interval lies below one.  In two dimensions we also
  display the joint histogram; the five-dimensional panel compares the same
  fixed critical projection and reports its independently calibrated product
  edges.

\subsection{Full sweep}

\begin{table}[H]
\centering
\small
\begin{tabular}{@{}ll@{}}
\toprule
Quantity & Values \\
\midrule
dimension $d$ & $2,5$ \\
block count $B$ (sample count $dB$) & $1,2,4,8$ \\
matrix law & broad-gap $(0.4,0.18)$; clean-mode $(0.2,0.10)$ \\
relative gap position & five points from $0.1$ to $0.9$ \\
mean nonlinear strength $\bar\beta$ & $0.46,0.50,0.54,0.555,0.575,0.595,0.65$ \\
coefficient law & constant; i.i.d. $0.555\pm0.04$; persistent $0.555\pm0.04$ \\
flip probability $p_{\rm flip}$ & $0.002,0.00035$ for persistent laws \\
label-noise scale $\sigma_\varepsilon$ & $0.005$ primary; $10^{-3}$--$10^{-8}$ in cusp audit \\
independent chains per main distribution setting & $192$ or $256$ \\
independent chains per reliability replicate & $96$ \\
steps per chain & $16{,}000$--$20{,}000$ \\
\bottomrule
\end{tabular}
\caption{Finite-step experiment grid.  The named matrix-law pair in
parentheses is $(\delta_\lambda,\phi)$.  Product diagnostics are separately
estimated with normalized products and a Feynman--Kac particle system.  Exact
numerical values and deterministic seed rules are included in the bundle.}
\label{tab:full-grid}
\end{table}

The experiments are structured sweeps rather than a claim of exhaustive
coverage: each panel changes one factor while holding the others fixed.  Null
results and stopped escape controls remain in the exported CSV/JSON files.

\section{Finite-step results}
\label{sec:results}

\subsection{Calibrated phase sweep}

For the broad-gap law in two dimensions, the full-product estimates give
$(\widehat\eta_{\mathrm L},\widehat\eta_{\mathrm H})=(2.128,2.268)$
at $B=1$.  Batch averaging moves both edges toward $2$ and narrows their
separation from $0.140$ at $B=1$ to $0.076$, $0.040$, and $0.019$ at
$B=2,4,8$, respectively.  Table~\ref{tab:threshold-results} reports the
local-interpolation Monte Carlo intervals.  These are product estimates: the harmonic curve is
computed with the projective particle system, not with a one-step norm.

\begin{figure}[H]
  \centering
  \includegraphics[width=\linewidth]{output/discrete_selfstab/figure1_discrete_product_thresholds.pdf}
  \caption{\textbf{Full-product threshold calibration.}
  Top Lyapunov exponent (left) and projective inverse-moment pressure
  $\log k_\eta(-1)$ (right) for $B=1,2,4,8$ Gram blocks.  Vertical markers
  are the higher-replication local crossing estimates; shaded regions are
  their 95\% Monte Carlo intervals.  Increasing $B$ narrows, but does not
  eliminate over the tested values $B\in\{1,2,4,8\}$, the
  post-Lyapunov/pre-harmonic interval.}
  \label{fig:discrete-phase-sweep}
\end{figure}

\subsection{Inside the Lyapunov--harmonic gap}

This is the primary experiment.  Across the broad-gap sweep, the median
radius increases from $0.043$ just after the Lyapunov edge to $0.139$ just
before the harmonic edge, while the $0.99$ quantile increases from $0.691$
to $1.135$.  The empirically estimated nonlinear tangent exponent remains
negative ($-0.100$ to $-0.167$ over those endpoints), even though the origin's linear
exponent is positive.  This is the measured finite-step self-stabilization
mechanism: expanding small states activate a nonlinear return at finite
amplitude.  All primary chains are finite and none is clipped or reset.

Across all structured sweeps and reliability replicas, the publication run
contains $35$ primary configurations and $129{,}760{,}000$ exact SGD
transitions.  It recorded zero nonfinite primary transitions; the largest
primary norm was $2.976$.  The separate five-noise-scale cusp audit adds
$76{,}800{,}000$ exact updates, bringing the total discrete numerical
evidence to $206{,}560{,}000$ finite-step transitions, again with zero
nonfinite iterates.

\begin{figure}[H]
  \centering
  \includegraphics[width=\linewidth]{output/discrete_selfstab/figure2_between_edges_distributions.pdf}
  \caption{\textbf{What happens between the two discrete thresholds.}
  Joint $(u_1,u_2)$ densities (top) and log-radius densities (bottom) below
  the Lyapunov edge, at three locations spanning the gap, and beyond the
  harmonic edge.  The broad-gap law becomes increasingly intermittent.  Its
  joint heatmap does not form the clean separated lobes of
  Figure~\ref{fig:discrete-bimodality}, although its critical-coordinate
  marginal develops a shallow central depression at the two post-harmonic
  points.  Every panel is an exact finite-step SGD chain.}
  \label{fig:gap-distributions}
\end{figure}

The negative full-product root in \eqref{eq:negative-root} is also estimated
at five positions in the gap.  Figure~\ref{fig:gap-scaling} is a deliberately
low-noise mechanism probe, not a measurement made in the primary
$\sigma_\varepsilon=0.005$ chains.  It compares the root with an empirical
small-amplitude exponent at $\sigma_\varepsilon=10^{-8}$.  The comparison is
labelled \emph{intermediate scaling}: fixed label noise rounds the density at
the origin, and the saturating feedback cuts off large excursions.  The
product roots at relative gap positions
$0.1,0.3,0.5,0.7,0.9$ are $0.105,0.321,0.534,0.736,0.897$; the corresponding
empirical small-ball slopes are $0.158,0.326,0.508,0.671,0.788$.  Repeating the
same fit on only the first half of each retained trajectory gives
$0.158,0.327,0.504,0.662,0.784$; the largest half/full-horizon change is
below $0.010$.  This nested check provides no evidence, at these two
horizons, that the discrepancy is driven by horizon length.  Agreement is
closest through the central part of the interval; finite noise and the
restricted fitting window remain limitations of this intermediate-scale
comparison.

\begin{figure}[H]
  \centering
  \includegraphics[width=\linewidth]{output/discrete_selfstab/figure8_gap_scaling.pdf}
  \caption{\textbf{Detailed scaling inside the Lyapunov--harmonic gap.}
  The full-product negative root and empirical intermediate small-ball exponents
  from the full and half retained horizons at bounded label-noise scale
  $10^{-8}$,
  small-amplitude cumulative probabilities, tail amplification and central
  occupancy, and nonlinear activation/tangent contraction at five relative
  gap positions.  This is a low-noise mechanism probe, not evidence that the
  cusp is visible at the primary noise scale, and it is not an asymptotic
  far-tail law.}
  \label{fig:gap-scaling}
\end{figure}

\subsection{Noise-scale and horizon audit of the cusp}
\label{sec:cusp-audit}

We next test the two-sided window principle in
Remark~\ref{rem:reinjection-window} rather than selecting a visually linear
segment after seeing the data.  At the authoritative broad-law
midpoint $\eta=2.198293$ and $\bar\beta=0.595$, the full-product prediction is
$\alpha=0.5284$.  For each noise scale, define the exact one-step standard
deviation of the update noise in the pinned coordinate by
\begin{equation}
 s_\xi=\eta\sigma_\varepsilon
 \sqrt{\E[(\Hvec_t)_{11}]/(dB)},
 \label{eq:cusp-innovation-scale}
\end{equation}
with $B=1$ in this audit, and define $r_{\rm cut}=0.092057$ by
$\bar\beta r_{\rm cut}^2/(1+r_{\rm cut}^2)=0.005$.  The fitting window is
predeclared as
\begin{equation}
 \max\{10s_\xi,\widehat q_{.01}\}<|u_1|
 <\min\{r_{\rm cut},\widehat q_{.50}\},
 \label{eq:cusp-fit-window}
\end{equation}
and is accepted only if the physical span
$r_{\rm cut}/(10s_\xi)$ is at least $10$ and the empirical fitting span is at
least $4$.  These guards are prespecified operational proxies for
$r_{\rm core}\ll|u_1|\ll r_{\rm nl}$; they do not verify the uniform renewal,
angular, or regeneration hypotheses of
Proposition~\ref{prop:matrix-reinjection}.

At the primary scale $\sigma_\varepsilon=0.005$, the physical span is only
$1.186$, so there is no admissible cusp window and we report no exponent.
The scale $0.001$ also fails, with span $5.929$; its superficially linear fit
is displayed as a rejected hollow marker.  Valid windows first appear at
$10^{-4}$, $10^{-5}$, and $10^{-8}$, with physical spans $59.29$, $592.86$,
and $592864$.  Their full-horizon fitted small-ball exponents, reported as
mean plus or minus standard deviation across three independent seeds, are respectively
$0.5706\pm0.0044$, $0.5629\pm0.0071$, and
$0.5309\pm0.0042$.  Thus the estimate approaches the product prediction
$0.5284$ as the intermediate window opens.

Each noise scale is run with $128$ chains per seed; the horizons
$12{,}000$, $24{,}000$, and $40{,}000$ are nested prefixes.  At $10^{-8}$ the corresponding
mean slopes are $0.5358$, $0.5336$, and $0.5309$, and the fit $R^2$ rises
from $0.99933$ to $0.99980$.  The low-noise agreement therefore survives the
reported horizon and seed checks.  More importantly, the audit rejects a
cusp claim at the paper's primary noise level.  It is consistent with the
renewal theorem's scale-separation requirement and the conditional matrix
extension, but does not verify the latter's angular or regeneration
assumptions.

\begin{figure}[H]
  \centering
  \includegraphics[width=\linewidth]{output/discrete_selfstab/figure9_cusp_noise_horizon_audit.pdf}
  \caption{\textbf{When is the product-root cusp observable?}
  Exact finite-step SGD at the broad-law midpoint across five label-noise
  scales, three independent seeds, and three nested horizons.  Fits use only the
  predeclared window \eqref{eq:cusp-fit-window}.  No admissible window exists
  at the primary $\sigma_\varepsilon=0.005$ or at $0.001$.  Once sufficient
  scale separation opens, the fitted small-ball exponent approaches the full-product
  prediction $\alpha=0.528$, reaching $0.531\pm0.004$ at $10^{-8}$ and
  remaining stable across the tested horizons.}
  \label{fig:cusp-noise-horizon-audit}
\end{figure}

\subsection{Deterministic nonlinear strength}

Since $\nabla^2V_\beta(\mathbf0)=0$, the independently estimated linear
diagnostics are identical across values of $\bar\beta$.  The nonlinear
distributions need not be.  We measure how the cutoff radius, burst duration,
return time, and possible mode locations depend on $\bar\beta$ at fixed
relative step position.

In a separate nonlinear-strength sweep at $\eta=2.195$, approximately $48\%$
of the way across the calibrated broad-law gap, increasing $\bar\beta$ from
$0.46$ to $0.595$
reduces the empirical $0.99$ radius from $1.111$ to $0.918$ while leaving the
linear thresholds unchanged.  Every tested far-field product exponent and
retained-horizon nonlinear tangent exponent is negative.  Thus the nonlinear
coefficient controls the cutoff and excursion size, not the location of the
inner product thresholds.

\begin{figure}[H]
  \centering
  \includegraphics[width=\linewidth]{output/discrete_selfstab/figure3_feedback_strength_sweep.pdf}
  \caption{\textbf{Deterministic self-stabilization strength.}
  At fixed matrix law and mid-gap step size, norm survival curves, median and
  $0.99$ radius, and the far-field/nonlinear tangent exponents are shown as
  $\bar\beta$ varies.  Stronger feedback lowers the nonlinear cutoff while
  the independently estimated linear diagnostics remain fixed.}
  \label{fig:beta-strength}
\end{figure}

\subsection{Stochastic and slowly varying coefficients}

The coefficient processes in \eqref{eq:beta-processes} keep the mean fixed
while varying amplitude and temporal persistence.  We test whether coefficient
variance or persistence changes excursion size, clustered bursting, return
times, and visible modes.  No universal tail exponent is imposed.

At $\eta=2.247$, approximately $85\%$ of the way across the calibrated gap,
and at matched mean $0.555$ and amplitude $0.04$, neither i.i.d. nor persistent
two-state modulation materially changes the radial tail amplification or
mean burst duration: $q_{.99}/q_{.50}$ stays between $8.81$ and $8.94$, while
mean burst duration stays between $2.239$ and $2.245$ steps.  Persistence does
increase the magnitude of the instantaneous coefficient--radius correlation
from about $0.014$ for i.i.d. modulation to $0.037$ and $0.046$ for
$p_{\rm flip}=0.002$ and $0.00035$.  All simulated augmented-chain
trajectories remained finite over the reported horizon.  Thus, within this
four-law ablation at amplitude $0.04$, we detect no material change in the
marginal radial tail.  This is a matched-mean null observation, not evidence
of invariance over coefficient amplitudes or persistence laws.

\begin{figure}[H]
  \centering
  \includegraphics[width=\linewidth]{output/discrete_selfstab/figure4_stochastic_feedback.pdf}
  \caption{\textbf{Bounded stochastic and slowly varying quartic curvature.}
  All runs have the same mean coefficient $\bar\beta$.  Curves compare a
  constant coefficient, i.i.d. two-state coefficients, and persistent
  two-state coefficients.  The panels show the matched-mean log-radius
  density, radial survival, tail amplification, burst duration, and
  coefficient--radius correlation.  The sampled coefficient is used once in
  the same finite SGD step as its batch Hessian.}
  \label{fig:stochastic-beta}
\end{figure}

\subsection{Directional bimodality, dimension, and batch size}

The additional broad-law metric sharpens what the joint heatmaps alone can
show.  The critical-coordinate center-to-side-mode ratio remains above one
just before the harmonic edge, $1.012$ $[1.007,1.017]$, then falls to
$0.957$ $[0.941,0.972]$ and $0.793$ $[0.779,0.806]$ at the two tested
post-harmonic points.  These are statistically resolved but shallow
directional depressions, not the clean separated joint lobes seen below.
Under the more strongly pinned clean-mode law, the ratio is $1.010$ inside
the gap, with 95\% chain-bootstrap interval $[1.005,1.015]$, so that run is
not classified as split.  It falls to $0.096$ $[0.088,0.106]$ just beyond
the harmonic edge and $0.0039$ $[0.0028,0.0051]$ deeper beyond it.  Thus both
experimental laws place the measured directional change after the harmonic
crossing, while the strength and geometric visibility of the lobes depend
strongly on directional pinning.  This supports a calibrated empirical
marker; it does not turn the harmonic diagnostic alone into a general
necessary-and-sufficient theorem.

\begin{figure}[H]
  \centering
  \includegraphics[width=\linewidth]{output/discrete_selfstab/figure5_discrete_bimodality.pdf}
  \caption{\textbf{Directional mode formation in the clean-mode law.}
  Critical-coordinate histograms (top) and joint 2D densities (bottom) below
  the Lyapunov edge, in the gap, and at two post-harmonic step sizes.  Two
  lobes in the pooled retained-sample histograms appear only in the displayed
  post-harmonic runs.}
  \label{fig:discrete-bimodality}
\end{figure}

A separate dimension comparison is run at $\eta=2.144$, selected from the
two-dimensional calibration as $\widehat\eta_{\mathrm H}^{(2)}+0.080$ and
lying between the two displayed two-dimensional post-harmonic settings.  At
this common nominal step size, the center-to-mode point estimates are
$0.0119$ $[0.0094,0.0147]$ in $d=2$ and
$0.0035$ $[0.0026,0.0047]$ in $d=5$.  An independent five-dimensional
product calibration gives
$(\widehat\eta_{\mathrm L}^{(5)},\widehat\eta_{\mathrm H}^{(5)})
=(2.029,2.061)$, placing the common step $0.083$ beyond its estimated
five-dimensional harmonic edge.  The step was nevertheless selected from
the two-dimensional calibration, so the figure retains that provenance in
its label.  For the broad-gap product, increasing $B$ from $1$ to $8$
substantially narrows the calibrated two-dimensional gap.
Because these are exact regression batches, the $B$ comparison changes both
the variability of the averaged Hessian and the innovation covariance, which
scales as $1/B$ by \eqref{eq:bounded-regression-innovation}.  The matched-gap
panel is therefore a batch-size comparison, not an intervention that holds
additive-noise variance fixed while changing only the matrix product.

\begin{figure}[H]
  \centering
  \includegraphics[width=\linewidth]{output/discrete_selfstab/figure6_dimension_batch_robustness.pdf}
  \caption{\textbf{Dimension and batch-size robustness.}
  Gap width versus $B$ (left), matched-relative-position log-radius densities
  for $B=1$ and $B=4$ (center), and clean-mode critical-coordinate densities
  at a common nominal step size in $d=2$ and $d=5$ (right).  The $d=5$ curve
  is labelled by how the common step was selected; a separate product
  calibration confirms that it also lies beyond the estimated $d=5$
  harmonic edge.}
  \label{fig:dimension-batch}
\end{figure}

\subsection{Convergence and failure accounting}

Independent-seed reruns agree on the median and $0.99$ radius.  The ordered
telescoping decomposition
$\log(r_{\rm lin}/r_{\rm old})+\log(r_{\rm nl}/r_{\rm lin})+
\log(r_{\rm new}/r_{\rm nl})=\log(r_{\rm new}/r_{\rm old})$ holds to machine
precision.  The three components depend on this stated ordering; they are not
an order-independent additive decomposition of the vector update.
The untamed negative-quartic and affine controls are stopped at first escape rather
than reset or recycled; they are never used as stationary samples.
At the stress step size $\eta=2.4$, all $20$ affine controls crossed
$\|u_t\|=10^6$ within $53$--$132$ steps, all $20$ zero-initialized untamed
negative-quartic controls crossed within $345$--$8651$ steps, and all six
untamed controls initialized at radius $3$ crossed within $39$--$67$ steps.

\begin{figure}[H]
  \centering
  \includegraphics[width=0.96\linewidth]{output/discrete_selfstab/figure7_discrete_audits.pdf}
  \caption{\textbf{Numerical reliability of the discrete chain.}
  Exact conditional log-radius balance; median and $q_{.99}$ across
  independent seeds plus a finite-radius initialization check; and
  first-escape times for affine and untamed negative-quartic controls.  The primary
  saturated-loss runs have zero nonfinite transitions, clipping operations,
  or resets.}
  \label{fig:reliability}
\end{figure}

\section{Interpretation and theorem status}
\label{sec:status}

The linear quantities $\gamma_\eta$ and $h_\eta$ do not depend on
$\beta_t$, because the quartic correction has zero derivative at the origin.
This does not make the nonlinear stationary law independent of $\beta_t$.
The experiments test two separate statements: whether the linear crossings
remain fixed when the nonlinear coefficient law is changed, and how its mean,
variance, and persistence change the finite-step stationary behavior at the
same calibrated location.

Below $\eta_{\mathrm L}$, Theorem~\ref{thm:kesten-baseline} supplies the
affine reference law under its assumptions.  Between $\eta_{\mathrm L}$ and
$\eta_{\mathrm H}$ we do not claim an upper-tail Kesten law.  Instead,
Theorem~\ref{thm:scalar-reinjection} proves a lower-amplitude on--off law for
an exact regenerative scalar excursion, while
Proposition~\ref{prop:matrix-reinjection} gives a conditional matrix result
under additional projective spectral, Markov-renewal, angular, and
regeneration assumptions.  At fixed
additive noise this is only an intermediate-window statement.  The primary
$\sigma_\varepsilon=0.005$ chains have no resolved cusp window; their direct
diagnostic is a recurrent cloud with intermittent excursions and a nonlinear
return force that becomes dominant at finite radius.  Beyond
$\eta_{\mathrm H}$ we report whether the critical projection actually becomes
bimodal.  Agreement supports a finite-step harmonic-bimodality hypothesis;
disagreement is not removed by redefining the threshold.

The Kesten statement below the Lyapunov edge and
Theorem~\ref{thm:ordering} are discrete-time results.  The two-step expansion
\eqref{eq:two-step-expansion} and Theorem~\ref{thm:exact-2d-flip} are also
exact for their stated maps.  For the nonlinear random matrix recursion,
however, neither $k_\eta(-1)=1$ nor the sign of a local quartic term alone is
a necessary-and-sufficient condition for bimodality.  Such a theorem would
require recurrence, reflection symmetry, one critical direction, regular
reinjection, and exclusion of competing global modes.  The scalar renewal
theorem proves why $\alpha=1$ changes the \emph{local density shape}; it does
not turn that local change into a theorem asserting exactly two global modes.
The present paper therefore treats the finite-step connection between
$h_\eta$ and modal structure as a tested hypothesis, while reporting every
exception.

For the particular bounded-Rademacher transition law, Feller continuity is
verified, whereas $\psi$-irreducibility, petiteness, and hence uniqueness and
positive Harris recurrence have not been established.  Likewise, the
negative-pressure calculations use the spectral identification in
Assumption~\ref{ass:negative-pressure-identification}.  These are explicit
analytical assumptions, not conclusions drawn from finite simulations.

\section{Conclusion}

Finite-step random-covariate SGD has three logically separate components: a
linear matrix product, a nonlinear return mechanism, and the geometry of the
resulting invariant law.  The top Lyapunov exponent locates loss of linear
stationarity.  For the bounded laws here, inverse-moment finiteness follows
from Lemma~\ref{lem:inverse-moment-finite}; under the spectral-identification
and strict-convexity assumptions of Theorem~\ref{thm:ordering}, an
inverse-moment crossing cannot precede the Lyapunov edge.  The exact
regenerative theorem identifies what lies between those two diagnostics:
the negative pressure root gives an integrable central cusp
$|a|^{\alpha-1}$ with $0<\alpha<1$, provided a finite-mean nonlinear
reinjection and a genuine additive-core/nonlinear-cutoff window exist.  Its
matrix extension additionally requires projective mixing and angular
pinning.  The exact
two-step flip map explains how a negative local quartic can
nevertheless stabilize alternating iterates, while the saturating potential
makes the global experiment well posed without altering the update after the
fact.  Calibrated experiments throughout the interval between the two linear
diagnostics determine whether the recurrent cloud remains centered, broadens,
or separates along a pinned direction.  Bounded random and persistent
quartic coefficients then reveal which aspects of self-stabilization depend
on the mean feedback and which depend on its temporal variability.  A
separate scale audit prevents overinterpretation: the cusp is unresolved at
the primary noise level and approaches the pressure prediction only after a
low-noise intermediate window opens.

\appendix

\section{Reproducibility checklist}

The result bundle contains:
\begin{enumerate}[leftmargin=2em]
  \item the exact source implementing \eqref{eq:exact-discrete-update};
  \item complete configuration files and random seeds;
  \item aggregate retained-sample summaries and every stopped escape control;
  \item deterministic seed rules for reproducing
  $\widehat\gamma_\eta$ and $\widehat k_\eta(q)$;
  \item the particle counts, product lengths, and burn-ins used by the
  pressure estimator;
  \item pooled mode ratios and all plotted density data;
  \item the five-noise-scale cusp-window audit, its predeclared acceptance
  rule, and seed-by-horizon slope table;
  \item a manifest mapping every table entry and figure panel to raw files.
\end{enumerate}

\section{Additional result tables}

\begin{table}[ht]
\centering
\scriptsize
\begin{tabular}{@{}cclccc@{}}
\toprule
$d$ & $B$ & matrix law & $\widehat\eta_{\mathrm L}$
& $\widehat\eta_{\mathrm H}$ & gap width \\
\midrule
$2$ & $1$ & broad-gap & $2.1282\ [2.1272,2.1292]$ & $2.2684\ [2.2676,2.2692]$ & $0.1401$ \\
$2$ & $2$ & broad-gap & $2.0635\ [2.0628,2.0642]$ & $2.1393\ [2.1368,2.1416]$ & $0.0758$ \\
$2$ & $4$ & broad-gap & $2.0318\ [2.0313,2.0323]$ & $2.0721\ [2.0693,2.0748]$ & $0.0403$ \\
$2$ & $8$ & broad-gap & $2.0181\ [2.0178,2.0184]$ & $2.0374\ [2.0352,2.0399]$ & $0.0193$ \\
$2$ & $1$ & clean-mode & $2.0290\ [2.0285,2.0295]$ & $2.0642\ [2.0616,2.0675]$ & $0.0352$ \\
$5$ & $1$ & clean-mode & $2.0289\ [2.0284,2.0293]$ & $2.0613\ [2.0597,2.0636]$ & $0.0325$ \\
\bottomrule
\end{tabular}
\caption{Discrete product thresholds.  Each bracket is a 95\% Monte Carlo
interval conditional on the local interpolation bracket: complete product
trajectories are resampled for $\eta_{\mathrm L}$ and independent full-SMC
replicates for $\eta_{\mathrm H}$.}
\label{tab:threshold-results}
\end{table}

\begin{table}[ht]
\centering
\scriptsize
\begin{tabular}{@{}llcccccc@{}}
\toprule
matrix law & regime & $\eta$ & $\bar\beta$
& nonfinite & median $r$ & $q_{.99}(r)$ & center/mode (95\% CI) \\
\midrule
broad & pre-L & $2.073$ & $0.595$ & $0$ & $0.025$ & $0.493$ & $1.101\ [1.099,1.103]$ \\
broad & just post-L & $2.131$ & $0.595$ & $0$ & $0.043$ & $0.691$ & $1.077\ [1.074,1.081]$ \\
broad & mid-gap & $2.198$ & $0.595$ & $0$ & $0.084$ & $0.931$ & $1.046\ [1.042,1.050]$ \\
broad & just pre-H & $2.266$ & $0.595$ & $0$ & $0.139$ & $1.134$ & $1.012\ [1.007,1.017]$ \\
broad & post-H & $2.323$ & $0.595$ & $0$ & $0.189$ & $1.286$ & $0.957\ [0.941,0.972]$ \\
broad & deep post-H & $2.373$ & $0.595$ & $0$ & $0.231$ & $1.421$ & $0.793\ [0.779,0.806]$ \\
clean & in gap & $2.047$ & $0.650$ & $0$ & $0.075$ & $0.373$ & $1.010\ [1.005,1.015]$ \\
clean & post-H & $2.109$ & $0.650$ & $0$ & $0.202$ & $0.542$ & $0.096\ [0.088,0.106]$ \\
clean & deep post-H & $2.159$ & $0.650$ & $0$ & $0.272$ & $0.636$ & $0.0039\ [0.0028,0.0051]$ \\
\bottomrule
\end{tabular}
\caption{Selected nonlinear finite-step outcomes.  ``Center/mode'' is the
smoothed density at zero divided by its positive-side maximum; values well
below one indicate two directional lobes.  Intervals resample complete chains;
all estimates use 301 bins and Gaussian smoothing of 1.5 bins.}
\label{tab:nonlinear-results}
\end{table}

\clearpage
\begin{thebibliography}{99}

\bibitem{Kesten1973}
H. Kesten.
\newblock Random difference equations and renewal theory for products of
random matrices.
\newblock \emph{Acta Mathematica}, 131:207--248, 1973.
\newblock \url{https://doi.org/10.1007/BF02392040}.

\bibitem{Goldie1991}
C. M. Goldie.
\newblock Implicit renewal theory and tails of solutions of random equations.
\newblock \emph{Annals of Applied Probability}, 1(1):126--166, 1991.

\bibitem{GuivarchLePage2012}
Z. Gao, Y. Guivarc'h, and E. Le Page.
\newblock Stable laws and spectral gap properties for affine random walks.
\newblock \emph{Annales de l'Institut Henri Poincar\'e, Probabilit\'es et
Statistiques}, 51(1):319--348, 2015.
\newblock \url{https://arxiv.org/abs/1108.3146}.

\bibitem{WangSteinsaltz2020}
F. Wang and D. Steinsaltz.
\newblock On transfer operators for Markovian products of invertible random
matrices.
\newblock 2020.
\newblock \url{https://arxiv.org/abs/2001.06865}.

\bibitem{DelMoral2004}
P. Del Moral.
\newblock \emph{Feynman--Kac Formulae: Genealogical and Interacting Particle
Systems with Applications}.
\newblock Springer, 2004.

\bibitem{MeynTweedie2009}
S. Meyn and R. Tweedie.
\newblock \emph{Markov Chains and Stochastic Stability}.
\newblock Cambridge University Press, second edition, 2009.

\bibitem{BuraczewskiDamekMikosch2016}
D. Buraczewski, E. Damek, and T. Mikosch.
\newblock \emph{Stochastic Models with Power-Law Tails: The Equation
$X=AX+B$}.
\newblock Springer, 2016.
\newblock \url{https://doi.org/10.1007/978-3-319-29679-1}.

\bibitem{GurbuzbalabanSimsekliZhu2021}
M. G\"urb\"uzbalaban, U. \c{S}im\c{s}ekli, and L. Zhu.
\newblock The heavy-tail phenomenon in SGD.
\newblock In \emph{Proceedings of the 38th International Conference on
Machine Learning}, volume 139 of PMLR, pages 3964--3975, 2021.
\newblock \url{https://proceedings.mlr.press/v139/gurbuzbalaban21a.html}.

\bibitem{HodgkinsonMahoney2021}
L. Hodgkinson and M. Mahoney.
\newblock Multiplicative noise and heavy tails in stochastic optimization.
\newblock In \emph{Proceedings of the 38th International Conference on
Machine Learning}, volume 139 of PMLR, pages 4262--4274, 2021.
\newblock \url{https://proceedings.mlr.press/v139/hodgkinson21a.html}.

\bibitem{CohenEtAl2021}
J. M. Cohen, S. Kaur, Y. Li, J. Z. Kolter, and A. Talwalkar.
\newblock Gradient descent on neural networks typically occurs at the edge
of stability.
\newblock In \emph{International Conference on Learning Representations},
2021.
\newblock \url{https://arxiv.org/abs/2103.00065}.

\bibitem{AhnZhangSra2022}
K. Ahn, J. Zhang, and S. Sra.
\newblock Understanding the unstable convergence of gradient descent.
\newblock In \emph{Proceedings of the 39th International Conference on
Machine Learning}, volume 162 of PMLR, pages 247--257, 2022.
\newblock \url{https://proceedings.mlr.press/v162/ahn22a.html}.

\bibitem{AroraLiPanigrahi2022}
S. Arora, Z. Li, and A. Panigrahi.
\newblock Understanding gradient descent on the edge of stability in deep
learning.
\newblock In \emph{Proceedings of the 39th International Conference on
Machine Learning}, volume 162 of PMLR, pages 948--1024, 2022.
\newblock \url{https://proceedings.mlr.press/v162/arora22a.html}.

\bibitem{DamianNichaniLee2023}
A. Damian, E. Nichani, and J. D. Lee.
\newblock Self-stabilization: The implicit bias of gradient descent at the
edge of stability.
\newblock In \emph{International Conference on Learning Representations},
2023.
\newblock \url{https://arxiv.org/abs/2209.15594}.

\bibitem{Arnold1998}
L. Arnold.
\newblock \emph{Random Dynamical Systems}.
\newblock Springer Monographs in Mathematics. Springer, 1998.
\newblock \url{https://doi.org/10.1007/978-3-662-12878-7}.

\bibitem{PlattSpiegelTresser1993}
N. Platt, E. A. Spiegel, and C. Tresser.
\newblock On--off intermittency: A mechanism for bursting.
\newblock \emph{Physical Review Letters}, 70(3):279--282, 1993.
\newblock \url{https://doi.org/10.1103/PhysRevLett.70.279}.

\bibitem{HeagyPlattHammel1994}
J. F. Heagy, N. Platt, and S. M. Hammel.
\newblock Characterization of on--off intermittency.
\newblock \emph{Physical Review E}, 49(2):1140--1150, 1994.
\newblock \url{https://doi.org/10.1103/PhysRevE.49.1140}.

\bibitem{OttSommerer1994}
E. Ott and J. C. Sommerer.
\newblock Blowout bifurcations: The occurrence of riddled basins and
on--off intermittency.
\newblock \emph{Physics Letters A}, 188(1):39--47, 1994.
\newblock \url{https://doi.org/10.1016/0375-9601(94)90114-7}.

\bibitem{Asmussen2003}
S. Asmussen.
\newblock \emph{Applied Probability and Queues}.
\newblock Stochastic Modelling and Applied Probability, volume 51.
Springer, New York, second edition, 2003.
\newblock \url{https://doi.org/10.1007/b97236}.

\bibitem{Kesten1974}
H. Kesten.
\newblock Renewal theory for functionals of a Markov chain with general
state space.
\newblock \emph{The Annals of Probability}, 2(3):355--386, 1974.
\newblock \url{https://doi.org/10.1214/aop/1176996654}.

\bibitem{NeyNummelin1987I}
P. Ney and E. Nummelin.
\newblock Markov additive processes I: Eigenvalue properties and limit
theorems.
\newblock \emph{The Annals of Probability}, 15(2):561--592, 1987.
\newblock \url{https://doi.org/10.1214/aop/1176992159}.

\bibitem{NeyNummelin1987II}
P. Ney and E. Nummelin.
\newblock Markov additive processes II: Large deviations.
\newblock \emph{The Annals of Probability}, 15(2):593--609, 1987.
\newblock \url{https://doi.org/10.1214/aop/1176992160}.

\bibitem{Alsmeyer1994}
G. Alsmeyer.
\newblock On the Markov renewal theorem.
\newblock \emph{Stochastic Processes and their Applications}, 50(1):37--56,
1994.
\newblock \url{https://doi.org/10.1016/0304-4149(94)90146-5}.

\end{thebibliography}

\end{document}
